\subsection{Calibrated intervals for the predictable residual}
The prediction target is the \emph{record-level segment-mean} off-dipole residual of the target
lead (one scalar per record: the mean over the segment's samples). We train a gradient-boosted
quantile regressor of it on the observed leads' segment means and apply Mondrian conformalized
quantile regression \cite{romano2019cqr} per $(S,s,\ell)$ group, with strict fold discipline:
basis and model on folds 1--7, the regressor's \texttt{max\_iter} selected by pinball loss on
fold 8, conformal calibration on fold 9, and a single evaluation on fold 10 (nominal
$1{-}\alpha{=}0.90$). This is a sound \emph{application} of established conformal machinery, not
a new predictor; the contribution is attaching it, per feature and lead, to the recoverability
map. Across the pre-registered $(S,s,\ell)$ groups, fold-10 within-group
coverage clusters near nominal ($0.87$--$0.94$; e.g.\ $\{$I,II,V1,V3,V5$\}$ QRS/V2
$\CalCov$, record-bootstrap CI $[\CalCovLo,\CalCovHi]$, mean width $\CalWidth$\,mV).

\textbf{What is and is not guaranteed.} Mondrian CQR gives \emph{within-$(S,s,\ell)$-group
marginal} coverage under exchangeability---\emph{not} coverage conditional on diagnosis, which
is not a conditioning variable. Diagnostic subgroups are therefore not guaranteed and the
smallest ones fall below nominal: the worst is the \WorstSubName{} subgroup at coverage
$\WorstSubCov$ ($n{=}\WorstSubN$), the expected finite-sample behaviour when conditioning on
a variable the calibration does not stratify on. We report worst-subgroup coverage rather than
hide it under the aggregate band, and we make no distribution-free claim beyond exchangeability
with the calibration fold.